% ------------------------------------------------------------------
% closure_section.tex -- new material for "The Link" (working draft, Aug 2026)
%
% Splice target: end of Section 5, REPLACING the boxed remark "Conservation
% of rents" (this proposition absorbs and upgrades it).
% Assumes the house preamble (amsthm environments: proposition / remark, as
% in the Unpriced-Value-of-Time paper); adjust environment names at splice
% time if the Link source differs. Every displayed claim is verified in
% checks/check_closure.py; the Prop 4(ii) patch at the bottom is verified in
% checks/check_baseline_props.py.
% ------------------------------------------------------------------

\subsection*{Closing the loop: where the spending goes}

Proposition 3 prices machines from the cost side; nothing so far says where
household spending comes to rest. This subsection closes the corner-regime
economy and turns ``conservation of rents'' from a bookkeeping remark into an
accounting identity.

\textbf{Setup.} Households jointly hold the land (the distribution of claims
is arbitrary here, and becomes the subject of the feasibility proposition to
follow). Preferences are Cobb--Douglas over the machine-made final good $g$
and direct land services $h$ --- housing, siting, space --- with land share
$\sigma \in (0,1)$: $U = g^{1-\sigma}h^{\sigma}$. Land is homogeneous with
total stock $T$, renting at $r$, split between production use $T_P$ (entering
machine-making through $\ell$) and housing $T_H$. Machines perform all tasks
(labor has exited, or is of negligible measure); $\delta = 0$, so the interest
sliver is suppressed.\footnote{With $\delta > 0$ a thin interest flow
reappears on the income side --- exactly the wage of waiting of
Proposition 3; nothing below changes beyond carrying the term.} Free entry in
machine-making gives $c = \ell r/(1-a)$ and hence
$p_g = c\bar\rho\bar L$.

\begin{proposition}[Closure and conservation of rents]
In the corner regime: (i) equilibrium exists and the real allocation is
unique,
\[
T_H = \sigma T, \qquad T_P = (1-\sigma)T, \qquad
Y \;=\; \frac{(1-a)(1-\sigma)\,T}{\ell\,\bar\rho\,\bar L};
\]
(ii) the entire goods bill is production-land rent, and national income is
aggregate site rent identically,
\[
p_g Y \;=\; r\,T_P, \qquad\qquad p_g Y + r\,T_H \;=\; rT;
\]
(iii) market clearing reproduces the cost-side relative price of
Proposition~4(ii), $r/p_g = (1-a)/(\ell\bar\rho\bar L)$: the model is
overidentified, and consistently so.
\end{proposition}

\begin{proof}
Gross machine services $X$ solve $X = \bar\rho\bar L\,Y + aX$ --- services for
final tasks plus services consumed making services --- so
$X = \bar\rho\bar L\,Y/(1-a)$, and production land demand is $\ell X$.
Substituting $p_g = \ell r\bar\rho\bar L/(1-a)$ gives
$p_g Y = r\,\ell\bar\rho\bar L\,Y/(1-a) = r\,\ell X = r\,T_P$: goods spending
passes through the zero-profit machine sector and lands, all of it, on
production land. Land is the only primary factor employed, so household income
is $rT$. Cobb--Douglas shares give $rT_H = \sigma rT$ and
$p_g Y = (1-\sigma)rT$, hence $T_H = \sigma T$; land-market clearing
$T_P = T - T_H$ then pins $Y$ as displayed; and dividing
$p_g Y = (1-\sigma)rT$ by the expression for $Y$ returns
$r/p_g = (1-a)/(\ell\bar\rho\bar L)$, which is (iii). Uniqueness is immediate
from the shares.
\end{proof}

Dissipation is migration, now by construction rather than assertion: every
unit of final spending resolves through the machine recursion into a claim on
a site. The escape route the informal version left open --- that automation's
released value might dissipate into consumer surplus at collapsing prices
rather than reattach to land --- is closed: prices do collapse \emph{and} the
identity holds, because what collapses is $p_g$ relative to $r$, never the
rent's claim on the flow of spending.

\begin{remark}[Substitution: the CES dial]
Replace Cobb--Douglas with CES elasticity $\eta$ between goods and land
services, and write $q \equiv r/p_g$. The land share of expenditure becomes
\[
s_h(q) \;=\; \frac{\sigma q^{1-\eta}}{\sigma q^{1-\eta} + 1 - \sigma},
\]
and as machine improvement sends $q \to \infty$ (Proposition~4(ii)):
$s_h \to 1$ for $\eta < 1$ --- complements; the squeeze sharpens and housing
eats the budget --- while $s_h \to 0$ for $\eta > 1$ and substitution bounds
it. Prediction~8 is thereby promoted from a deflator observation to an
estimate of $\eta$; the climbing shelter share of low-income budgets is the
$\eta < 1$ signature.
\end{remark}

\begin{remark}[Heterogeneous land and the idle margin]
With heterogeneous parcels the aggregates are unchanged: $rT$ reads
$\int r(\tau)\,d\tau$ over parcels in use, parcels at the idle margin earn
their reservation rent of zero, and $r$ in the identities is the rent
schedule, not a single number. Nothing in the proposition requires every acre
to be employed --- only that spending, wherever it goes, crosses a site.
\end{remark}

\begin{remark}[Corner-above]
With labor still present at machine parity the identity weakens to:
\emph{all non-wage income is site rent},
$\mathrm{GDP} = N c\bar\rho + rT$, with the wage bill pinned at parity and
vanishing along the $\bar\rho \to 0$ margin, recovering the identity of the
text in the limit. The full mixed-case accounting adds notation without
changing where anything goes.
\end{remark}

\begin{remark}[Distribution, made stark]
The identity prices the stakes of Proposition~6: in the corner, a household
without land claims has zero market income --- not low; zero. Whether rents
can fund the floor is therefore not one policy question among several; it is
the question of whether anyone without a site eats. The feasibility
proposition takes it up next.
\end{remark}

% ------------------------------------------------------------------
% PATCH -- erratum to Proposition 4(ii).
% (Caught by checks/check_baseline_props.py: the a -> 1 limit of
%  r/pg = (1-a)/(l rho Lbar) is ZERO, not infinity.)
%
% Replace the clause
%   "and it diverges along every margin of machine improvement:
%    rho -> 0, a -> 1, or l -> 0"
% with the following:
% ------------------------------------------------------------------
% ... and it diverges as the task-side edge closes ($\bar\rho \to 0$) or the
% recipe's land content vanishes ($\ell \to 0$). The self-replication margin
% moves it boundedly, and in the opposite direction to a naive reading:
% improvement is $a \downarrow$, which raises $r/p_g$ toward its bound
% $1/(\ell\bar\rho\bar L)$, while $a \to 1$ is regress --- it sends machine
% costs, not rents, to infinity ($c = \ell r/(1-a)$), and $r/p_g \to 0$.
% Even substituting machines for land inside the recipe,
% $\ell = \ell_0(1-a)$, leaves $r/p_g = 1/(\ell_0\bar\rho\bar L)$: bounded.
% ------------------------------------------------------------------
